\chapter{问题背景、重述与分析}

\section{问题背景}
自动化车床在离散制造过程中承担连续加工任务，刀具状态直接决定产品质量与产线效率。若刀具进入失效状态而未及时识别，将持续产出不合格品并触发返工或报废成本；若检查过于频繁或提前换刀，则会增加检查费用与刀具消耗费用。该问题本质上是“故障风险控制”与“预防性维护投入”之间的平衡优化问题。

题目给出的背景具有三个典型特征：
\begin{enumerate}
  \item \textbf{随机性}：刀具寿命存在显著波动，故障出现时刻不可预知；
  \item \textbf{信息不完备}：生产过程无法直接观测刀具状态，只能通过抽检产品间接判断；
  \item \textbf{成本耦合}：检查、误判、故障恢复、提前换刀等多类成本相互影响。
\end{enumerate}

\section{题目重述}
已知：题目提供 100 次刀具故障记录（故障发生时累计加工件数）与费用参数
\[
\theta_1=200,\ \theta_2=10,\ \theta_3=3000,\ \theta_4=1000\quad (\text{单位：元}).
\]
在第二类场景中还给出误判相关信息：正常状态下不合格率为 $2\%$，故障状态下合格率为 $40\%$，且误停损失费用 $\theta_5=1500$ 元/次。

需解决以下三问：
\begin{enumerate}
  \item 在“正常全合格、故障全不合格”条件下，设计最优检查间隔与换刀策略；
  \item 在存在误判条件下，重新设计最优策略；
  \item 讨论改进检查方式（如双检确认）是否可进一步提升经济效益。
\end{enumerate}

\section{研究目标与技术路线}
本文目标是建立统一的期望成本模型，将策略设计转化为约束优化问题，并给出可复现的求解流程。技术路线如下：
\begin{enumerate}
  \item 对刀具寿命样本进行分布分析，得到概率模型；
  \item 对“检查节点 + 换刀决策”进行参数化表达；
  \item 计算各策略下的合格品期望、废品期望、误判期望和总成本期望；
  \item 以单位合格品平均成本最小为准则搜索最优策略；
  \item 对关键参数做灵敏度分析，评估策略稳定性。
\end{enumerate}

\section{主要符号说明}
\begin{table}[H]
\centering
\caption{核心符号定义}
\begin{tabular}{ll}
\toprule
符号 & 含义 \\
\midrule
$m$ & 刀具寿命（加工件数）随机变量 \\
$f(m)$ & 刀具寿命概率密度函数 \\
$k_1,\ldots,k_n$ & 检查节点对应的产品序号 \\
$R_1,R_2$ & 正常/故障状态下产品合格率 \\
$\zeta_1,\zeta_2,\zeta_3$ & 合格品、总产量、废品期望 \\
$\gamma$ & 检查次数期望 \\
$\eta_1,\eta_2$ & 误判停机期望、主动换刀概率 \\
$h$ & 一个生产周期总损失期望 \\
$\phi$ & 单位合格品平均成本 \\
\bottomrule
\end{tabular}
\end{table}

\section{本章小结}
本章明确了问题结构、目标函数构建方向和符号系统。后续章节将依次给出数学模型、求解算法与结果分析。
